\chapter{Lesson 1: The Sounds of Mandarin}

\section{Mandarin and pinyin}

Modern Standard Chinese (\zh{普通话} pǔtōnghuà, ``the common speech''; in Taiwan
\zht{國語} guóyǔ) is based on the pronunciation of Beijing. It is written with
characters (\zh{汉字} Hànzì), which do not reliably show pronunciation, so in the
first weeks of the course we learn \emph{pinyin} (\zh{拼音} pīnyīn, ``spell sound''),
the official romanization adopted in the PRC in 1958. Pinyin is a tool for learning and
for typing; Chinese people read characters, not pinyin.

\begin{definition}[The Mandarin syllable]
Every Mandarin syllable has three parts:
\begin{itemize}
  \item an \emph{initial} (\zh{声母} shēngmǔ): the consonant at the start (it may be
        missing);
  \item a \emph{final} (\zh{韵母} yùnmǔ): the vowel part, possibly ending in
        \emph{-n} or \emph{-ng};
  \item a \emph{tone} (\zh{声调} shēngdiào): the pitch contour of the whole syllable.
\end{itemize}
For example \emph{mǎ} = initial \emph{m} + final \emph{a} + third tone.
\end{definition}

There are only about 400 distinct initial--final combinations (about 1300 with tones),
far fewer than English syllables. That is why tones matter: they are what tells
\emph{mā} ``mother'' apart from \emph{mǎ} ``horse''.

\begin{center}
\begin{tikzpicture}[box/.style={draw, thick, rounded corners, minimum height=0.9cm, minimum width=1.6cm}]
  \node[box, fill=blue!10] (I) at (0,0) {\emph{zh}};
  \node[box, fill=orange!15, minimum width=2.2cm] (F) at (2.1,0) {\emph{ong}};
  \node[box, fill=green!15] (T) at (4.3,0) {1st tone};
  \node[font=\large] at (6.4,0) {$=$};
  \node[font=\large] at (7.8,0) {\emph{zhōng} \zh{中}};
  \node[below=0.15cm of I, font=\scriptsize] {initial};
  \node[below=0.15cm of F, font=\scriptsize] {final};
  \node[below=0.15cm of T, font=\scriptsize] {tone};
\end{tikzpicture}
\end{center}

\section{Initials}

The 21 initials are grouped by where in the mouth they are made. The teacher drilled
them in this order, which is also the order of the pinyin chart.

\begin{center}
\begin{tabular}{llll}
\toprule
Group & Initials & Place of articulation & Example \\
\midrule
Labials & b p m f & lips & \emph{bà} \zh{爸}, \emph{pà} \zh{怕}, \emph{mā} \zh{妈}, \emph{fàn} \zh{饭} \\
Alveolars & d t n l & tongue tip on gum ridge & \emph{dà} \zh{大}, \emph{tā} \zh{他}, \emph{nǐ} \zh{你}, \emph{lǎo} \zh{老} \\
Velars & g k h & back of tongue & \emph{gē} \zh{哥}, \emph{kàn} \zh{看}, \emph{hǎo} \zh{好} \\
Palatals & j q x & tongue flat behind lower teeth & \emph{jiā} \zh{家}, \emph{qù} \zh{去}, \emph{xiè} \zh{谢} \\
Retroflexes & zh ch sh r & tongue tip curled back & \emph{zhè} \zh{这}, \emph{chī} \zh{吃}, \emph{shì} \zh{是}, \emph{rén} \zh{人} \\
Dental sibilants & z c s & tongue tip behind upper teeth & \emph{zài} \zh{在}, \emph{cài} \zh{菜}, \emph{sì} \zh{四} \\
\bottomrule
\end{tabular}
\end{center}

\subsection{Aspiration, not voicing}

English speakers hear \emph{b} vs \emph{p} as voiced vs voiceless. In Mandarin both are
voiceless; the difference is \emph{aspiration} (a puff of air).

\begin{center}
\begin{tabular}{lll}
\toprule
Unaspirated & Aspirated & Test \\
\midrule
b & p & hold a tissue in front of your mouth: it should move only for \emph{p} \\
d & t & \\
g & k & \\
j & q & \\
zh & ch & \\
z & c & \\
\bottomrule
\end{tabular}
\end{center}

\begin{important}[Letters that do not sound like English]
\begin{itemize}
  \item \textbf{q} is like the \emph{ch} of ``cheese'' with the lips spread; \textbf{x} is
        between \emph{s} and \emph{sh}. So \emph{qī} \zh{七} is roughly ``chee'' and
        \emph{xiè} \zh{谢} roughly ``shyeh''.
  \item \textbf{c} is \emph{ts} as in ``cats'' (aspirated): \emph{cài} \zh{菜} ``tsai''.
  \item \textbf{z} is \emph{ds} as in ``kids'': \emph{zài} \zh{再}.
  \item \textbf{zh, ch, sh, r} curl the tongue back. \emph{r} is between English \emph{r}
        and the \emph{s} of ``measure''.
  \item \textbf{h} is rougher than English \emph{h}, closer to the \emph{ch} of Scottish
        ``loch''.
\end{itemize}
\end{important}

\section{Finals}

There are 36 or so finals. The chart groups them by their first vowel.

\begin{center}
\begin{tabular}{llll}
\toprule
Group & Finals & Example syllables & Notes \\
\midrule
Simple & a o e i u ü & \emph{bā, mó, hē, nǐ, bù, lǜ} & \emph{e} is a back unrounded ``uh'' \\
Compound & ai ei ao ou & \emph{lái, gěi, hǎo, dōu} & \emph{ao} like ``ow'' in ``how'' \\
Nasal (-n) & an en in un ün & \emph{kàn, mén, nín, lùn, jūn} & \emph{un} = \emph{uen} \\
Nasal (-ng) & ang eng ing ong & \emph{máng, děng, míng, zhōng} & \\
i-group & ia ie iao iu ian iang iong & \emph{jiā, xiè, xiǎo, liù, tiān, liǎng, xióng} & \emph{ian} sounds ``yen'' \\
u-group & ua uo uai ui uan uang & \emph{huā, duō, kuài, duì, huān, huáng} & \\
ü-group & üe üan & \emph{xué, yuǎn} & \\
Retroflex & er & \emph{èr} \zh{二} & \\
\bottomrule
\end{tabular}
\end{center}

\subsection{The special vowel after z, c, s, zh, ch, sh, r}

After these seven initials the letter \emph{i} does not sound like ``ee''. It is a
buzzing continuation of the consonant: \emph{sì} \zh{四} ``four'' is roughly
``sz'', and \emph{shì} \zh{是} ``to be'' is roughly ``shr''. Compare \emph{xī}
(``shee'') with \emph{shī} (``shr''): mixing these up is the most common mistake in
the first quizzes.

\section{Tones}

Mandarin has four tones plus a \emph{neutral tone}. Linguists describe pitch on a scale
from 1 (lowest in your range) to 5 (highest).

\begin{center}
\begin{tikzpicture}[x=1cm, y=0.55cm]
  \foreach \y in {1,...,5} {
    \draw[gray!40] (0,\y) -- (11.2,\y);
    \node[left, font=\scriptsize] at (0,\y) {\y};
  }
  \draw[->] (0,0.5) -- (0,5.8) node[above, font=\scriptsize] {pitch};
  % tone 1: 55
  \draw[very thick, blue] (0.4,5) -- (2.0,5);
  \node[font=\small] at (1.2,0) {1st: 55};
  \node[font=\small] at (1.2,-0.9) {\emph{mā} \zh{妈}};
  % tone 2: 35
  \draw[very thick, red!70!black, -{Stealth}] (2.8,3) -- (4.4,5);
  \node[font=\small] at (3.6,0) {2nd: 35};
  \node[font=\small] at (3.6,-0.9) {\emph{má} \zh{麻}};
  % tone 3: 214
  \draw[very thick, green!50!black, -{Stealth}] (5.2,2) .. controls (5.8,0.6) and (6.3,0.6) .. (6.8,4);
  \node[font=\small] at (6.0,0) {3rd: 214};
  \node[font=\small] at (6.0,-0.9) {\emph{mǎ} \zh{马}};
  % tone 4: 51
  \draw[very thick, orange!80!black, -{Stealth}] (7.6,5) -- (9.2,1);
  \node[font=\small] at (8.4,0) {4th: 51};
  \node[font=\small] at (8.4,-0.9) {\emph{mà} \zh{骂}};
  % neutral
  \fill[purple] (10.4,2) circle (0.12);
  \node[font=\small] at (10.4,0) {neutral};
  \node[font=\small] at (10.4,-0.9) {\emph{ma} \zh{吗}};
\end{tikzpicture}
\end{center}

\begin{center}
\begin{tabular}{lllll}
\toprule
Tone & Mark & Description & Example & Meaning \\
\midrule
1st & \emph{ā} & high and level, as if singing a note & \zh{妈} \emph{mā} & mother \\
2nd & \emph{á} & rising, like a surprised ``what?'' & \zh{麻} \emph{má} & hemp; numb \\
3rd & \emph{ǎ} & low dipping, then rising & \zh{马} \emph{mǎ} & horse \\
4th & \emph{à} & sharply falling, like ``No!'' & \zh{骂} \emph{mà} & to scold \\
neutral & \emph{a} & short and light, no mark & \zh{吗} \emph{ma} & question particle \\
\bottomrule
\end{tabular}
\end{center}

\begin{aside}
Mnemonic from class: \emph{mā mà mǎ}---``Mom scolds the horse''. The professor also had
us say the four tones of \emph{ma} while moving a hand: flat, up, down-and-up, down.
In connected speech the third tone almost never rises fully; it stays low (the ``half
third tone'', pitch 21). The full 214 is heard mainly on its own or at the end of a
sentence. See Lesson 2.
\end{aside}

\subsection{The neutral tone}

Some syllables lose their tone and become short and light. Their pitch depends on the
preceding tone (high after the 3rd tone, low after the 4th). Neutral tones occur in
\begin{itemize}
  \item grammatical particles: \zh{吗} ma, \zh{呢} ne, \zh{吧} ba, \zh{的} de, \zh{了} le;
  \item the second syllable of reduplicated kin terms: \zh{妈妈} māma, \zh{爸爸} bàba,
        \zh{哥哥} gēge;
  \item many common words: \zh{谢谢} xièxie, \zh{先生} xiānsheng, \zh{朋友} péngyou.
\end{itemize}

\section{Spelling rules}

Pinyin has a few rules that change how a final is written. They do not change the sound.

\begin{important}[Pinyin spelling conventions]
\begin{enumerate}
  \item \textbf{y and w.} A syllable with no initial that begins with \emph{i}, \emph{u}
        or \emph{ü} is written with \emph{y} or \emph{w}: \emph{i} $\rightarrow$ \emph{yi},
        \emph{ia} $\rightarrow$ \emph{ya}, \emph{in} $\rightarrow$ \emph{yin};
        \emph{u} $\rightarrow$ \emph{wu}, \emph{uo} $\rightarrow$ \emph{wo};
        \emph{ü} $\rightarrow$ \emph{yu}, \emph{üe} $\rightarrow$ \emph{yue}.
  \item \textbf{ü without dots.} After \emph{j, q, x} and \emph{y} the dots are dropped,
        because \emph{u} cannot follow these letters anyway: \emph{jù, qù, xué, yú}. After
        \emph{n} and \emph{l} the dots stay, since \emph{nu} and \emph{nü} are both
        possible: \emph{nǚ} \zh{女} ``female'' vs \emph{nù} \zh{怒} ``anger''.
  \item \textbf{Shortened finals.} With an initial, \emph{iou, uei, uen} are written
        \emph{iu, ui, un}: \emph{liù} \zh{六}, \emph{duì} \zh{对}, \emph{lùn} \zh{论}.
  \item \textbf{Apostrophe.} When a syllable starting with \emph{a, o, e} follows another
        syllable, an apostrophe marks the boundary: \emph{Xī'ān} \zh{西安} (two syllables)
        vs \emph{xiān} \zh{先} (one).
\end{enumerate}
\end{important}

\subsection{Where does the tone mark go?}

\begin{enumerate}
  \item If there is an \emph{a} or \emph{e}, it takes the mark: \emph{hǎo, xiè, guān}.
  \item In \emph{ou}, the \emph{o} takes it: \emph{dōu, zǒu}.
  \item Otherwise the \emph{last} vowel takes it: \emph{liù, duì, guó}.
\end{enumerate}
When marking \emph{i}, the dot is replaced by the mark (\emph{nǐ}).

\section{Classroom expressions}

These were on the first vocabulary quiz.

\begin{center}
\begin{tabular}{lll}
\toprule
Characters & Pinyin & Meaning \\
\midrule
\zh{你好} & nǐ hǎo & hello \\
\zh{老师好} & lǎoshī hǎo & hello, teacher \\
\zh{同学们好} & tóngxuémen hǎo & hello, class \\
\zh{谢谢} & xièxie & thank you \\
\zh{不客气} & bú kèqi & you're welcome \\
\zh{对不起} & duìbuqǐ & sorry \\
\zh{没关系} & méi guānxi & it doesn't matter \\
\zh{再见} & zàijiàn & goodbye \\
\zh{请问} & qǐngwèn & may I ask \ldots \\
\zh{请再说一遍} & qǐng zài shuō yí biàn & please say it again \\
\zh{我不懂} & wǒ bù dǒng & I don't understand \\
\zh{请跟我说} & qǐng gēn wǒ shuō & please repeat after me \\
\zh{上课} / \zh{下课} & shàngkè / xiàkè & class begins / class is over \\
\bottomrule
\end{tabular}
\end{center}

\begin{aside}
Numbers 1--10 were taught in week one because they are perfect tone practice:
\zh{一} yī, \zh{二} èr, \zh{三} sān, \zh{四} sì, \zh{五} wǔ, \zh{六} liù, \zh{七} qī,
\zh{八} bā, \zh{九} jiǔ, \zh{十} shí. Note how \emph{sì} (4) and \emph{shí} (10) differ in
both initial and tone: a classic listening test.
\end{aside}

\section*{Review}
\begin{itemize}
  \item A syllable is initial + final + tone; tones distinguish meaning.
  \item Aspiration (not voicing) separates \emph{b/p, d/t, g/k, j/q, zh/ch, z/c}.
  \item \emph{q, x, c, z} and the retroflex series are the hardest consonants;
        after \emph{z, c, s, zh, ch, sh, r} the letter \emph{i} is a buzz, not ``ee''.
  \item Tones: 55, 35, 214, 51, plus a light neutral tone on particles and many second
        syllables.
  \item Spelling: \emph{y/w} at the start of a syllable, \emph{ü} written \emph{u} after
        \emph{j, q, x, y}, \emph{iu/ui/un} are shortened, apostrophe before \emph{a/o/e}.
  \item Tone marks go on \emph{a} or \emph{e}, then the \emph{o} of \emph{ou}, otherwise
        the last vowel.
\end{itemize}
